%% ============================================================================
%% IV. Hypotheses                                             target: ~0.4 page
%% Stage-1 reviewers score (i) the value of the hypotheses and (ii) whether the
%% planned evaluation can support OR falsify them. Make each hypothesis
%% testable and state what result would refute it. Empirical hypotheses are OK.
%% Reference them as \ref{hyp:...} (prints H1, H2, ...).
%% ============================================================================
\section{Hypotheses}
\label{sec:hypotheses}

\todo{Refine the wording and add quantitative thresholds once the setup is fixed.}

\begin{hypothesis}[Reachability]\label{hyp:reach}
Allowing a single contact between the rod and an environment edge strictly
enlarges the set of reachable tip poses, in particular the range of tip
directions, compared with contact-free manipulation in the same constrained
workspace ($\Rfree\subsetneq\Rcon$).
\end{hypothesis}
\begin{outline}
  \item Rationale: the contact acts as an ad-hoc joint (pivot), and its
    location along the rod is controlled by feeding the rod over the edge.
  \item Refuted if, for the tested clearances, the tip directions reachable
    with contact do not exceed those reachable without it (simulation and
    hardware).
\end{outline}

\begin{hypothesis}[Vision-free estimation]\label{hyp:estimate}
Gripper pose, wrist wrench, and contact location suffice to estimate the tip
pose without vision, and the contact location significantly reduces the error
compared with gripper pose and wrench alone.
\end{hypothesis}
\begin{outline}
  \item Rationale: without contact, the wrench at the gripped end determines
    the rod's shape~\cite{bretl2014quasistatic,mishani2022realtime}; a contact
    adds an unknown reaction force, and knowing where it acts recovers the
    missing information.
  \item Refuted if the tip error on unseen configurations exceeds the task
    tolerance, or if adding the contact location gives no significant
    improvement (input ablation, \cref{sec:experiments}).
\end{outline}

\begin{hypothesis}[Contact-aided steering]\label{hyp:steer}
Closed-loop steering with the haptic tip-pose estimate brings the tip to
target poses behind the edge with a higher success rate than contact-free and
open-loop strategies.
\end{hypothesis}
\begin{outline}
  \item Refuted if the success rate and final error are not better than those
    of the baselines on the same targets.
\end{outline}

\begin{outline}
  \item \textbf{Expected insights:} when contact helps and by how much (as a
    function of rod stiffness, clearance, and contact location), and which
    haptic signals carry the tip-pose information.
\end{outline}
